\scanpage{53}
résultant de l'autre suite spectrale pour la cohomologie de De Rham) ne peut guère être autre chose que l'homomorphisme de Frobenius. Il faut dire que tout ça est bien éculé, et qu'on ne pourra dire \deleted{\unsure{dans|des}} des choses vraiment intéressantes et nouvelles qu'en faisant appel à la ``\unsure{vraie|vcaie}'' cohomologie p\unsure{-}adique.

Bien cordialement
\margin{52}
